% ===== §16 =====
\section{Non-Possession as the Condition Throughout}
\label{sec:nonpossession}

\noindent
Non-possession has appeared at each turn of the model, in the restraint that lets the thing thrown remain the other's, in the withholding that is non-intervention, in the trust that preserves without demanding a return. It is time to state its place directly. Non-possession is not an eighth capacity, one more power alongside the seven. It is the condition under which the seven are what they are, and without which each of them declines into a form of the extraction the normative case set against understanding. It runs through the model as a single requirement in seven forms, and the model is completed by naming it as such.

A word on the standing of this condition holds it to the fuller ground drawn in the normative case. The legitimacy of understanding was there observed across several ethical traditions, of which non-possession is one aspect, nearest the refusal of the other as a mere means and completed by the attentive response, the reciprocity, and the fruitfulness the other traditions render. Non-possession serves as the condition of the seven capacities because it is the aspect that governs, most directly, the conduct of the crossing, what one does and refrains from doing in the exercise of each capacity. It carries the fuller legitimacy into the practice of the model without standing in for the whole of it, and the attentiveness, the mutuality, and the fruitfulness that the other traditions render are present in the capacities as the positive content that non-possession, a restraint, leaves for them to supply. The condition named here is thus the practical face, in the exercise of the capacities, of the legitimacy observed across the traditions.

The point can be seen by taking the seven capacities and removing the condition from each, so that what remains is the capacity's counterfeit. The recognition of the other's world, without non-possession, becomes the reconnaissance of a territory to be taken, the acknowledgment of the other's difference as the first move in mapping it for a hold. The entry into the other's world becomes an infiltration, an entry made to secure a position and not to observe with the other. The structural mapping becomes the diagram by which the other is managed, the translation a means of putting the other's words to one's own uses. The throwing of a generative thing becomes the planting of a suggestion the other is led to think their own, the most refined of manipulations, which borrows the very form of the generative gift to install a conclusion by stealth. Non-intervention becomes the calculated withholding that lets the other commit themselves before one moves. And the preservation of relational experience becomes the keeping of an account, a record of what the other owes, held against a future reckoning. Each capacity has its counterfeit, and in every case the counterfeit is the capacity emptied of non-possession and turned to a hold.

This is why non-possession cannot be one capacity among the others, for it is the condition that separates each of the others from its counterfeit, and a model that listed it as a further item would misplace it as one more thing to do, where it is the standing under which all the others are done. The same power that enters the other's world can enter to observe with the other or to secure a hold upon them, and nothing in the capacity itself decides which, for the capacity is the same power in both cases and its direction is given by the condition under which it is exercised. Non-possession is that condition, the orientation of the whole crossing toward observing with the other and not toward possessing them, and it is present or absent in each capacity as the capacity is exercised, not as a separate act performed before or after.

The wrong that the counterfeits share is the one the normative case named, and it is legible in the terms of generative justice. In each counterfeit the generativity of the relation is turned to the account of one party, the other made into a resource from which value is drawn under the appearance of a bond in which value is shared. The other becomes, in the phrase the wider framework uses, fuel for an end that is not the relation's own, and the more skillfully the counterfeit imitates the legitimate capacity the more thoroughly the extraction is concealed. Non-possession is the refusal of this turn, the holding of the relation's generativity as belonging to the relation and to both who generate within it, and the treatment of the other, throughout the crossing, as a co-generating subject whose world is entered to be observed with and never annexed.

It remains only to hold this condition, as the normative case held its criterion, within the history the account theorizes, so that non-possession is not installed as the timeless law the framework elsewhere refuses. What non-possession requires, where the line falls between a legitimate entry and an annexation, between a generative gift and a planted suggestion, between a preservation and an account kept, is not fixed for all times. It is drawn in the historical development of relational life and will be drawn again as that life develops, and the model offers non-possession as the form the condition takes in the present, held as the best the present affords and open to its own further generation. With this the model is complete, seven capacities in two orders of cultivation, under a condition that runs through all seven and that is itself, like everything in the account, generated in the history it belongs to and set as no final guarantee.
